\section*{Chapter \ref{ch:g6:wholes} --- Whole Numbers}

\begin{solution}{exo:g6:wholes:1}
$3\,047$; \quad $20\,500$; \quad $104$. \quad
$60\,013$: ``sixty thousand and thirteen''.
\end{solution}

\begin{solution}{exo:g6:wholes:2}
Hundreds digit: $2$; units digit: $9$; the $4$ counts the
ten-thousands. Decomposition:
\[
47\,209 = 4 \times 10\,000 + 7 \times 1000 + 2 \times 100 + 0 \times 10
+ 9 .
\]
\end{solution}

\begin{solution}{exo:g6:wholes:3}
$507 < 570$ (tens: $0 < 7$); \quad
$1\,099 > 989$ (four digits beat three); \quad
$23\,456 < 23\,465$ (tens: $5 < 6$); \quad
$9\,999 < 10\,000$ (four digits against five).
\end{solution}

\begin{solution}{exo:g6:wholes:4}
By digit count, $978$ and $1\,001$ come before the three
four-digit numbers starting with $2$; among those, compare the tens and
units: $2\,034 < 2\,043 < 2\,304$. Final order:
\[
978 < 1\,001 < 2\,034 < 2\,043 < 2\,304 .
\]
\end{solution}

\begin{solution}{exo:g6:wholes:5}
$348 + 576 = 924$; \quad $805 - 268 = 537$ (check:
$537 + 268 = 805$); \quad $307 \times 26 = 307 \times 20 + 307 \times 6
= 6\,140 + 1\,842 = 7\,982$.
\end{solution}

\begin{solution}{exo:g6:wholes:6}
$38 + 45 + 12 + 5 = (38 + 12) + (45 + 5) = 50 + 50 = 100$.

$2 \times 83 \times 50 = (2 \times 50) \times 83 = 100 \times 83 =
8\,300$.

$4 \times 78 \times 25 = (4 \times 25) \times 78 = 100 \times 78 =
7\,800$.
\end{solution}

\begin{solution}{exo:g6:wholes:7}
$95 = 7 \times 13 + 4$ (\omterm{def:g3:division:remainder}{quotient} $13$, \omterm{def:g3:division:remainder}{remainder} $4$).

$230 = 6 \times 38 + 2$ (\omterm{def:g3:division:remainder}{quotient} $38$, \omterm{def:g3:division:remainder}{remainder} $2$).

$504 = 8 \times 63 + 0$ (\omterm{def:g3:division:remainder}{quotient} $63$, \omterm{def:g3:division:remainder}{remainder} $0$: $504$ is
\omterm{def:g6:wholes:divisible}{divisible} by $8$).
\end{solution}

\begin{solution}{exo:g6:wholes:8}
The equality is true, but the \omterm{def:g3:division:remainder}{remainder} $13$ is \emph{not} smaller than
the \omterm{def:g5:division:multiple}{divisor} $8$, so it is not the Euclidean division. Take one more
$8$ out of the \omterm{def:g3:division:remainder}{remainder}: $173 = 8 \times 21 + 5$, with $0 \leq 5 < 8$:
\omterm{def:g3:division:remainder}{quotient} $21$, \omterm{def:g3:division:remainder}{remainder} $5$.
\end{solution}

\begin{solution}{exo:g6:wholes:9}
\emph{1.} $2\,000 = 12 \times 166 + 8$: the farm fills $166$ full
boxes and $8$ eggs are left over.

\emph{2.} $150 = 12 \times 12 + 6$: twelve boxes contain only $144$
eggs, not enough. The bakery must buy $13$ boxes.
\end{solution}

\begin{solution}{exo:g6:wholes:10}
$100 = 7 \times 14 + 2$: one hundred days are $14$ full weeks and $2$
days more. Fourteen weeks later it is again a Tuesday; two days after a
Tuesday is a \emph{Thursday}.
\end{solution}

\begin{solution}{exo:g6:wholes:11}
The largest possible \omterm{def:g3:division:remainder}{remainder} in a division by $9$ is $8$ (the
\omterm{def:g3:division:remainder}{remainder} must be smaller than the \omterm{def:g5:division:multiple}{divisor}). So the dividend is
\[
9 \times 56 + 8 = 504 + 8 = 512 .
\]
\end{solution}

\begin{solution}{exo:g6:wholes:12}
Let the digits be $h$ (hundreds), $0$ (tens), $u$ (units). We know
$h = 2 \times u$ and $h + 0 + u = 9$, so $2u + u = 9$, giving
$3u = 9$, $u = 3$, and $h = 6$. The number is $603$. Check: $6$ is
twice $3$, the tens digit is $0$, and $6 + 0 + 3 = 9$.
\end{solution}
